\chapter{Оптические часы}\label{ch:ch0}

Тулий -- 69-й элемент таблицы Менделеева, входит в группу лантаноидов.
В лантаноидах, элементах с 57 по 71, происходит заполнение $4f$ оболочки, которая расположенная внутри уже заполненных $5s$, $5p$ и $6s$ оболочек. 
Из-за этого, переходы внутри этой оболочки, как правило, экранированы от внешних статических электрических полей, защищены от столкновений и влияния теплового излучения окружающей среды.
К такому типу переходов относятся часовые переходы в атомах тулия и эрбия \cite{patscheider2021observation}.

Единственный стабильный изотоп атомов тулия --- $^{169}$Tm является бозоном, спин его ядра $I = 1/2$, из-за чего каждый электронный уровень расщеплен на две сверхтонкие компоненты.
Структура основного состояния атомов тулия $\ket{[Xe]4f^{13}(^2 F^o)6s^2, J = 7/2}$ содержит один неспаренный электрон, вследствие чего имеет большой магнитный момент в основном состоянии ($\mu = 4 \mu_B$).
Это, и характерное для лантаноидов большое количество резонансов Фешбаха \cite{kotochigova2014controlling}, позволяющих контролировать межатомное взаимодействие, делают ансамбли нейтральных атомов тулия перспективной платформой для разработки на его основе квантового симулятора \cite{davletov2020machine}.
Разработка такого устройства ведётся в Российском квантовом центре.

В метрологии интерес к атомам тулия возникает из-за наличия магнитодипольного перехода на длине волны $\lambda = 1.14\unit{мкм}$ с естественной шириной линии $\gamma = 1.2\unit{Гц}$, который связывает между собой два уровня тонкой структуры основного состояния. 
Исследование этого перехода началось ещё в 80-х годах прошлого века, когда было показано, что в парах буферных газов он не испытывает столкновительного уширения \cite{александров1983неуширяемая}. 
Разработка же оптических часов на основе ансамблей нейтральных атомов тулия ведётся с $2008 \unit{года}$ в ФИАНе.
За это время было показано, что малы сдвиги частоты часового перехода, связанные с тепловым излучением окружения, а также со столкновениями атомов и статическим электрическим полем \cite{golovizin2019inner}.
Многие другие сдвиги частоты подавляются стандартными методами: сдвиг, вызываемый линейным эффектом Зеемана подавляется с помощью опроса перехода между центральными магнитными подуровнями, а сдвиги связанные с эффектом Доплера и импульсом отдачи фотона подавлены из-за опроса часового перехода в оптической решётке, что обеспечивает режим Лэмба-Дике.
Подавлению сдвигов частоты, связанных с квадратичным эффектом Зеемана и тензорной поляризуемостью с помощью формирования т.н. ``синтетической частоты'' посвящена работа \cite{golovizin2021simultaneous}.
Вкратце, идея этого метода состоит в том, что для перехода $\ket{g, F = 3} \xrightarrow{} \ket{c, F = 2}$ эти сдвиги равны по модулю и противоположены по знаку таким сдвигам для перехода $\ket{g, F = 4} \xrightarrow{} \ket{c, F = 3}$. 
Возбуждая оба этих перехода одновременно, и вычисляя полусумму их частот можно получить величину, не испытывающую этого сдвига. 

На данный момент, наибольшая неопределённость частоты в тулиевых оптических часах связанна с излучением оптической решётки. 
Дело в том, что при достижении уровня неточности $10^{-17}$ требуется учёт не только сдвигов, вызываемых электродипольной поляризуемостью, но и сдвигов, связанных с поляризуемостями более старших порядков (электроквадрупольной, магнитодипольной, а также гиперполяризуемостью).
Сложность их характеризации заключается в том, что они зависят не только от характеристик оптической решётки (её глубины, и отстройки от магической длины волны), но и от движения атомов.
Данная работа посвящена подготовке к их измерению.
В этой главе приводится обзор текущего протокола опроса атомов тулия, и мотивируется разработка нового.

На данный момент создана уже два прототипа такого устройства: экспериментальная установка для спектроскопии нейтральных атомов тулия, а также прототип аналогичной, но компактной системы для транспортируемой версии оптических часов \cite{golovizin2021compact}.
Наличие двух установок позволяет исследовать сдвиги частоты часового перехода с помощью сравнения двух систем.
Данная глава описывает основные этапы работы с нейтральными атомами тулия в часовом эксперименте, при этом в тех местах где это важно, будет делаться акцент на различиях между двумя установками.
Принципиальная схема импульсов в одном цикле работы с атомами тулия приведена на рисунке \ref{fig:PulseScheme}.


\begin{figure}[ht]
	\centerfloat{
		\includegraphics[scale=1]{Chapter0/PulseScheme}
	}
	\caption{Схема импульсов применяемая для подготовки начальных состояний атомов тулия в начальных состояниях часовго перехода.
	}\label{fig:PulseScheme}
\end{figure} 

%Также, в конце этой главы приводится проект следующего шага в создании транспортируемого репера частоты, и описание необходимых для этого приготовлений.
%В этой главе будет описано состояние эксперимента по созданию тулиевых оптических часов на момент начала данной работы.

%\footnote Так как запись часового перехода в атомах тулия $\ket{[Xe]4f^{13}(^2 F^o)6s^2, J = 7/2} \xrightarrow{} \ket{[Xe]4f^{13}(^2 F^o)6s^2, J = 5/2}$ достаточно громоздка, поэтому, далее в тексте будем обозначать основное состояние $\ket{[Xe]4f^{13}(^2 F^o)6s^2, J = 7/2}$ как $\ket{g}$, а часовой уровень $\ket{[Xe]4f^{13}(^2 F^o)6s^2, J = 5/2}$, как $\ket{c}$, при необходимости указывая полные моменты атома и/или его проекцию на ось квантования (например $\ket{g, F = 4}$ или $\ket{c, F = 3, m_F = 0}$). Остальные используемые в работе уровни также получат подобные обозначения.

\section{Старый протокол}

\subsection{Захват атомов в МОЛ}

Для получения атомов, образец металлического тулия нагревается в атомной печке до температур порядка $800\unit{К}$, что позволяет получить достаточно большой для дальнейшей работы поток атомов.
Однако при такой температуре лишь малая их доля обладает достаточно низкими скоростями, чтобы позволить захват в магнито-оптическую ловушку. 
Стандартным методом решить эту проблему является торможение атомов в Зеемановском замедлителе.
Зеемановское замедление и первая стадия охлаждения основана на переходе с длиной волны \fst со спектральной шириной $\gamma = 10\unit{МГц}$, что позволяло захватить несколько миллионов атомов охлаждённых до температур порядка $T\sim100\unit{мкК}$, что всё ещё слишком много для эффективного перезахвата в оптическую решётку.
Вторая стадия лазерного охлаждения используется для получения ансамбля атомов, который может быть перезахвачен в оптическую решётку.
Она основана на переходе с длиной волны \snd  со спектральной шириной $\gamma = 350\unit{кГц}$, что позволяет охлаждать атомы до температур порядка $T\sim20\unit{мкК}$ при полном перезахвате атомов из первичной МОЛ.
Схема уровней, и переходов, упоминающихся в данной работе изображена на рисунке \ref{fig:Levels}.

При создании второго поколения установки делался акцент на компактизации вакуумной камеры.
%Существует несколько подходов для этого как то пирамидальные МОЛ и так далее. 
Отличие второго поколения установки от первого состоит в отсутствии Зеемановского замедлителя: атомная печка расположена всего в $11\unit{см}$ от области захвата атомов.
При наличии тормозящего пучка, вместо катушек замедлителя используется остаточное магнитное поле МОЛ и маленькая дополнительная катушка, которые вместе создают магнитное поле, позволяющее осуществить Зеемановское торможение. 
При меньшей эффективности торможения большое число атомов удаётся захватить из-за большего телесного угла из которого происходит захват атомов.  
Вторым отличием является геометрия завода охлаждающего излучения в область захвата атомов: в случае большой установки две пары пучков находятся в горизонтальной плоскости, третья же пара заводится под небольшим углом к вертикали.
В компактной же установке две пары пучков заводятся под углом $45^\circ$ к вертикали, третья пара заведена вдоль оси X.

\begin{figure}[ht]
	\centerfloat{
		\includegraphics[scale=0.75]{Chapter0/Levels.pdf}
	}
	\caption{
		Схема уровней и переходов, используемых в работе.
		Переходы ``1st stage'' и ``2nd stage'' используются для захвата и охлаждения атомов в магнито-оптической ловушке.
		Переход ``Pumping'' использовался ранее для подготовки атомов тулия, захваченных в оптическую решётку на центральных магнитных подуровнях обоих сверхтонких компонент основного состояния.
		Переходы ``Clock'' используются в качестве часовых.
		Переход ``3rd stage'' исследуется в данной работе. С его помощью будет осуществляться охлаждение атомов в МОЛ, а также на его основе будет реализован протокол подготовки начальных состояний в оптической решётке. 
		На длине волны  ``Lattice'' формируется оптическая решётка. 
	}\label{fig:Levels}
\end{figure} 

Измерение числа атомов производится по измерению сигнала флуоресценции атомного ансамбля КМОП камерой \textbf{BASLERNAME}.
Для этого резонансно возбуждается перход на длине волны \fst, при этом часть рассеянных фотонов регистрируется камерой. 
Полученная фотография облака атомов аппроксимируется двумерной функцией Гаусса, при этом из параметров функции можно получить размер облака и число атомов. 
Это измерение разрушающее, при регистрации атомы нагреваются достаточно, чтобы покинуть область захвата.
Для измерения количества атомов на верхнем сверхтонком подуровне основного состояния \gF{3} используем этот же, но с перекачивающим излучением, настроенным в резонанс с 3-4 переходом. 

\subsection{Оптическая решётка}

Принцип работы оптической решётки описан в приложении \ref{app:A}. 
Для опроса часового перехода в атомах тулия используется вертикальная одномерная оптическая решётка на близкой к магической длине волны $\lambda=1063.5\unit{нм}$ (энергия отдачи $\Erec = h\times1.043\unit{кГц}$).
Источником излучения на этой длине волны служит волоконный лазер \textbf{PreciLaserModelName}, с выходной мощностью до $10\unit{Вт}$.
Стабильность частоты этого лазера лучше, чем $10\unit{МГц/ч}$, чего более чем достаточно для наших экспериментов.
В этой работе частота лазера не стабилизировалась.

В обеих установках для формирования оптической решётки использовались усиливающие резонаторы, что позволяло, во-первых, получать большие глубины потенциала (характерная глубина оптической решётки в эксперименте составляла $U = 1000\unit{E}_{\text{rec}}$), и во-вторых, формировать известную и предсказуемую моду резонатора.
В компактной системе резонатор внутривакуумный, конфокальный, с резкостью $\mathrm{F} = 400$ и усилением $100$, что позволяло получить до $10\unit{Вт}$ циркулирующей в резонаторе мощности при входной на уровне $100\unit{мВт}$.
В большой же системе использовался внешний резонатор, что приводило к гораздо меньшей резкости $\mathrm{F} = 15$, и как следствие, меньшему усилению ($4$). 
Для управления мощностью резона тора используются акусто-оптические модуляторы, что позволяет быстро манипулировать глубиной решётки.

В получаемы таким образом оптические решётки удавалось перезахватить до $30\%$ атомов из магнито-оптических ловушек, после чего устанавливались задающие магнитные поля для дальнейшего эксперимента.

\subsection{Подготовка начальных состояний}

После перезахвата атомов из магнито-оптической ловушки они были распределены по магнитным подуровням основного состояния. 
Для опроса часового перехода было необходимо подготовить их на центральном магнитном подуровне. 
Для этого использовался метод оптической накачки линейно-поляризованным излучением на длине волны $418\unit{нм}$ \cite{fedorova2020simultaneous}.
Идея этого метода базируется на том факте, что при возбуждении перехода между подуровнями с равным значением полного момента, запрещённым оказывается переход между центральными магнитными подуровнями. 
Таким образом, нужное состояние оказывается тёмным, и в нём скапливаются атомы.
Для используемого перехода, сверхтонкие расщепления основного и возбуждённого состояния оказались близкими, что вместе со спонтанным распадом на верхнее сверхтонкое позволяло получить атомы на центральном магнитном подуровне обоих сверхтонких подуровней.
При этом, для переноса населённости с крайних магнитных состояний на центральное, требуется рассеяние десятков фотонов, в результате чего происходит нагрев атомов, и, как следствие, потери из оптической решётки. 
В итоге, на двух центральных магнитных подуровнях, суммарно остаётся около $40\%$ атомов, захваченных в оптическую решётку.  

\begin{figure}[ht]
	\centerfloat{
		\includegraphics[draft=True]{Chapter0/PolarizationCleaning}
	}
	\caption{
		Что-то про старую накачку и потери в процессе.
	}\label{fig:OldPumping}
\end{figure} 

Для перераспределения атомов между двумя начальными состояниями часовых переходов, а также для удаления атомов, не накаченных на центральный магнитный подуровень, используется микроволновый переход между ними на частоте $1496\unit{МГц}$ \cite{mishin2024combined}. 
Идея вычищения поляризации атомного ансамбля изображена на рисунке \ref{fig:PolarizationCleaning} и состоит в следующем: после переноса атомов на верхнюю сверхтонкую компоненту $\ket{g, F = 3}$, невозбуждённые атомы удаляются из области захвата с помощью возбуждения перехода на длине волны \snd.
Этот переход достаточно узкий, чтобы в процессе не затрагивать атомы на верхней сверхтонкой компоненте.
Для возбуждения этого перехода использовалась самодельная дипольная антенна, источником сигнала для которой служил генератор SRS SG382.

\begin{figure}[ht]
	\centerfloat{
		\includegraphics[scale=0.5]{Chapter0/PolarizationCleaning}
	}
	\caption{
		Идея вычищения поляризации атомов с помощью микроволнового перехода между центральными магнитными подуровнями уровней тонкой структуры основного состояния.
		Вначале атомы как-то распределены по магнитным подуровням, причём большая их часть находится на центральном магнитном подуровне.
		Возбуждается переход, и все атомы из состояния $\ket{g, F=4, m_F =0}$ переносятся на $\ket{g, F=3, m_F =0}$.
		Невозбуждённые атомы удаляются из области захвата, после чего происходит произвольное перераспределение атомов между центральными магнитными подуровнями. 
	}\label{fig:PolarizationCleaning}
\end{figure} 

Для понижения конечной температуры ансамбля атомов использовался метод просеивания оптической решётки \cite{mishin2022effect}. 
Он заключался в кратковременном понижении глубины удерживающего потенциала для удаления наиболее горячих атомов и последующей термализации ансамбля при меньшей температуре. 
Таким образом получалось достичь среднего колебательного числа ансамбля вплоть до $\bar{n}_z \approx 0.5$ при числе атомов порядка $10^3$.
Недостатком такого метода было то, что чем ниже желаемая температура, тем меньше атомов в итоге оставалось жить в ловушке, что ограничивало снизу получаемые температуры.
Также, такой метод оказывался не слишком стабильным: конечная температура довольно сильно зависела от начальных характеристик ансамбля, поэтому приходилось калибровать её каждый раз заново.
Характерный спектр боковых колебательных частот часового перехода для тех времён представлен на рисунке \ref{fig:SidebandsOld} (Подробнее про измерение температуры атомов по спектроскопии боковых колебательных частот в приложении \ref{app:A}).   


\begin{figure}[ht]
	\centerfloat{
		\includegraphics[scale=0.5]{Chapter3/3_2_НастоящиеСайдбэнды}
	}
	\caption{
	}\label{fig:SidebandsOld}
\end{figure} 

\subsection{Часовой переход и синтетическая частота}

Источником излучения на длине волны $1.14\unit{мкм}$ служил лазер TopticaDL Pro стабилизированный по внешнему высокодобротному резонатору методом Паунда-Древера-Холла. 
Излучение лазера делилось на две системы, в каждой из них использовались два двухпроводных АОМа для возбуждения 43 и 32 переходов соответственно.
Далее, излучения совмещались друг с другом и доставлялись в систему по оптоволокну, где совмещались с излучением оптической решётки. 
Для измерения числа возбуждённых атомов использовался возврат на основные состояния коротким $\pi$-импульсом \cite{provorchenko2020clock}.

\subsection{Результаты}

С одной стороны, представленный выше метод был довольно прост в реализации, что несомненно было его плюсом. 
С другой стороны, мы испытывали значительные проблемы при попытке измерить поляризуемости высоких порядком в таком протоколе: температура атомов оказывалась достаточно высокой, нестабильной, и могла меняться со временем, что опять таки, плохо сказывалось на стабильности устройства.
Более того, получить атомы на конкретном колебательном подуровне тоже не получалось.
Всё это привело к необходимости разработки нового протокола подготовки атомов к опросу часового перехода.

\section{Разработка нового протокола}

Выделим несколько мест, которые могли бы быть улучшены в представленном выше протоколе.
%
Во-первых, более эффективное охлаждение атомов в магнито-оптической ловушке могло бы одновременно увеличить эффективность перезахвата в оптическую решётку и уменьшить дальнейшие потери атомов засчёт меньшей начальной температуры захваченного ансамбля.
Значительно улучшить протокол охлаждения атомов с помощью использованных ранее переходов, по видимому невозможно: подробнее в работах \cite{вишнякова2016вторичное, фёдорова2019исследование}.
Можно использовать для охлаждения другой охлаждающий переход, c меньшей спектральной шириной, чем у используемых (доплеровский предел температуры $T_D = \frac{h\gamma}{2k_B}$).
Переходы нужного типа, приведены в таблице из базы данных NIST \cite{NIST_ASD} (табл. \ref{tab:Transitions}).
Среди представленных в этой таблице переходов наименьшей спектральной шириной обладает переход на длине волны $506.2\unit{нм}$.
Вторая глава этой работы посвящена созданию магнито-оптической ловушки атомов тулия на длине волны \thd.

\begin{table}[]
	\centering
	\begin{tabular}{|c|c|r|c|c|c|r|}
		\cline{1-3} \cline{5-7}
		№ & $\lambda$, нм & $\gamma$, кГц & \multirow{11}{*}{   } &№ & $\lambda$, нм & $\gamma$, кГц\\
		\cline{1-3} \cline{5-7}
		1 & 251.4 & 843.5  && 10 & 356.7 & 604.8 \\
		\cline{1-3} \cline{5-7}
		2 & 255.0 & 429.7  && 11 & 371.8 & 22918.3 \\
		\cline{1-3} \cline{5-7}
		3 & 259.6 & 2228.2  && 12 & 375.2 & 276.9 \\
		\cline{1-3} \cline{5-7}
		4 & 262.2 & 827.6 && 13 & 410.6 & 10185.9 \\
		\cline{1-3} \cline{5-7}
		5 & 265.1 & 60.5 && 14 & 420.4 & 386.7 \\
		\cline{1-3} \cline{5-7}
		6 & 266.0 & 73.2 && 15 & 459.9 & 82.7 \\
		\cline{1-3} \cline{5-7}
		7 & 318.0 & 165.5 && 16 & 506.2 & 7.8 \\
		\cline{1-3} \cline{5-7}
		8 & 323.4 & 682.2 && 17 & 530.7 & 345.3 \\
		\cline{1-3} \cline{5-7}
		9 & 341.0 & 1559.7 && 18 & 567.6 & 206.9 \\
		\cline{1-3} \cline{5-7}
	\end{tabular}
	\caption{Переходы с основного состояния атомов тулия, которые могут быть использованы как охлаждающие.}
	\label{tab:Transitions}
\end{table}

Во-вторых, можно было бы изменить протокол подготовки атомов тулия в оптической решётке, включив в него охлаждение до основного колебательного состояния, для чего также подходит выбранный переход.
Подробнее о методах в главе \ref{ch:ch3}, которая будет посвящена разработке такого протокола.

Помимо охлаждения до основного колебательного состояния, для измерения поляризуемостей высоких порядков также необходимо уметь подготавливать ансамбли атомов тулия в различных колебательных состояниях. 
Разработка метода, позволяющего это сделать также будет представлена в главе \ref{ch:ch3}.